%*******************************************************
% Acknowledgments
%*******************************************************
%\pdfbookmark[1]{Acknowledgments}{acknowledgments}
\chapter*{Acknowledgments}
\addcontentsline{toc}{chapter}{Acknowledgments}
\chaptermark{Acknowledgments}
\begin{flushright}{\slshape
   .} \\ \medskip
    ---
\end{flushright}


As this thesis reaches it final pages, it is time to finally thank and acknowledge the people that have been part of this journey, in one way or the other. 

First of all, I want to thank Rudi Grimm. He took a bit of a bet by taking a lost theoretician, with almost no practical experience, to work in one of his labs, but he trusted me (more than I probably trusted myself) with this opportunity and I am more than happy to have done my PhD in the Dy-K experiment. Here I learned more than I ever thought I could. I got to learn from his wisdom, his knowledge of our field, and his experience in all academia-related matters.


Next is the Dy-K Team. I thank Elisa Soave and Zhu-Xiong Ye, for teaching me the Dy-K machine and experimental physics more in general, and for the hours spent in the lab scanning magnetic fields and looking at weak molecular signals. Yaakov Yudkin for his work and great documentation on the digital PID, and his availability while learning the VHDL code for the FPGA and while debugging the boards. 

Thanks to Marian Kreyer, for being a great research companion and friend over all the years of this PhD. From discussing the simulation code and his thesis, to the endless discussions (physics and non-physics-related), to the detailed feedback on my thesis, to the climbing and to the mountains related tips. If thesis would have second author Marian would be the one for this. 

Thanks to Chun-Kit Wong (Jacky) and Luc Absil, the new generation of the experiment. Together we fought endless battles against the force of evil, namely the ion pump of the dysprosium oven and all of the other vacuum related issues, and with them the Dy(-K) experiment keeps advancing. I  also want to acknowledge the master's students and interns that worked and orbited our Lab during the years: Nicola Grani,  Viviana Lippolis, Charlotte Farcot-Delfini, and Francesco Galasso.

A particular mention goes to Emil Kirilov, the fire brigade of our experiment, always ready to be the heavy muscle in our fight with the vacuum, and always ready to discuss any technical aspect of our experiments and to find the best ``Bulgarian'' way to solve our practical issues.

I also want to extend my thanks to the rest of the Grimm group: Cosetta Baroni, Erich Dobler, Isabella Fritsche, Adrian K\"onig, Bo Huang, Ruben Erlenstedt, Johanna Hennebichler, Gregor Anich, Niclas H\"ollrigl, for stimulating discussions and silly chats during lunch breaks.

A group like ours could not work without a great administrative support. I want to thank first of all Verena Tanzer,for her efficiency and her constant willingness to help, and for being the real engine of our department. I want also to extend my gratitude to the IQOQI administrative team and to the electronic and mechanical workshop of the university for their help over the years.


During my PhD I had the privilege to be part of the Doktoratskolleg atoms, light and molecules (DK-ALM), which has been a great opportunity for me. The DK-ALM provides students with stronger support from the university compared to the standard phd course of study, and it's my strong belief that every phd student should benefit from being part of it. For the first half of my studies the DK-ALM also financially supported me and allowed me to spend periods of time abroad. I thank Heide Streicher and Eric Endres for their help, organization of amazing Science days and Summer Schools, and their never ending patience with my delays in reply to their emails.


 
My first trip ``abroad'' was in Florence, in the Li-Cr experiment of Matteo Zaccanti. Despite being a relative short stay of few weeks I felt immediately part of the team, and I had the opportunity of learning a lot during my time there. Moreover, the knowledge exchange between our groups over the years is a great example of doing science in my opinion. Thanks to the whole team: Matteo Zaccanti, Adreas Trenkwalder, Alessio Ciamei, Stefano Finelli and Maximilian Schemmer. In particular I want to thank Matteo, which, together with the time I spent in Florence, is in no small part a reason why I could muster motivation and inspiration to complete my PhD, in a particularly low period of my scientific path. Additionally, I also want extend a special thank to Stefano, our path during the PhD have been quite parallel and sharing similar challenges in our restricted field of mass-imbalanced Fermi mixtures.

My second research stay was in Paris. I thank Sylvain Nascimbene, Jean Dalibard, Nehal Mittal, Tristan Villain, Evgenii Gadylshin, and all the other members of the BEC team at the Collège de France for the hospitality, and opportunity to learn and discover a different use of dysprosium, compared to ones typically available in Innsbruck. 


A PhD doesn't only mean 

 

mulitple phd friends  (name a few: Sarah, dessi, Kristoph, Thomas)
ferlainos Stv



Among all these wonderful people a special additional mention for YAO2026


additional climbing and mountain people


Conference people: Les Houches,


Making a step back in time, I would like to take advantage of this space to extend my gratitude to the people with whom I shared some part of my life with me during my time at the university in Trento, \graffito{Procrastinators of the world, unite!\\ \tiny{\dots tomorrow}}in part because I didn't have time to write any acknowledgments for my Master thesis at the time and in part because if I was well-equipped to finish my PhD is also thanks to them.

First of all I want to express my gratitude to my physics teachers. Starting from Iacopo Carusotto, my master 
Trento
physics teachers
Iacopo,  Luca, Giacomo, Alessio

via prepositura 

cibofisico

il Sanba

Gente di brescia





